%% ma: L12-B12-0021
%% lop: 12
%% chuong: 4
%% bai: 12
%% dang: 12.12.13
%% loai: DS
%% muc: [TH, VD, TH, VD]
%% dap_an: "SĐĐS"
%% nguon: "(NBV)-ĐỀ SỐ 18-MỖI NGÀY 1 ĐỀ THI 2025.pdf (D18, MỖI NGÀY 1 ĐỀ - Đề số 18), Phần II Câu 3"
%% kien_thuc: "Bài 11 (nguyên hàm), Bài 12 (tích phân, quãng đường từ vận tốc), đổi đơn vị m/s sang km/h"
%% trang_thai: chinh-thuc
%% ghi_chu: "Đáp án gốc SĐĐS đúng (giải độc lập: v = 0,0025t^2; l = 180 m; v(80) = 16 m/s = 57,6 km/h; thời gian qua cầu 660/16 = 41,25 s). Đề gốc có hình minh họa tàu, cây, cầu chỉ là tranh trang trí (không có số liệu toán) nên bỏ hình. Đã duyệt 19/10/2026."
\begin{ex}
Một đoàn tàu đang đứng yên trong sân ga, ngay trước đầu tàu có một cái cây. Đoàn tàu khởi hành từ trạng thái đứng yên với gia tốc $a=0{,}005t$ $(\mathrm{m/s^2})$ và đi qua cái cây trong thời gian $60$ giây. Sau $80$ giây đoàn tàu chuyển sang trạng thái chuyển động đều.
\choiceTF
{Vận tốc của đoàn tàu là $v=5\cdot10^{-3}t^2$ $(\mathrm{m/s})$.}
{\True Chiều dài của đoàn tàu là $l=180$ $(\mathrm m)$.}
{\True Sau $80$ giây, đoàn tàu chuyển động với tốc độ $57{,}6$ $(\mathrm{km/h})$.}
{Sau khi chuyển động đều một thời gian, đoàn tàu gặp một cây cầu có chiều dài $480$ $(\mathrm m)$. Khi đó đoàn tàu đi qua cây cầu đó trong thời gian $30$ giây.}
\loigiai{a) $v(t)=\displaystyle\int0{,}005t\,\mathrm dt=0{,}0025t^2+C$, mà $v(0)=0$ nên $v(t)=2{,}5\cdot10^{-3}t^2$, khác $5\cdot10^{-3}t^2$. Sai.

b) Tàu đi qua cái cây trong $60$ giây nên chiều dài tàu bằng quãng đường đi được trong $60$ giây: $l=\displaystyle\int_0^{60}0{,}0025t^2\,\mathrm dt=\dfrac{0{,}0025\cdot60^3}{3}=180$ (m). Đúng.

c) $v(80)=0{,}0025\cdot80^2=16$ (m/s) $=16\cdot3{,}6=57{,}6$ (km/h). Đúng.

d) Tàu chuyển động đều với vận tốc $16$ m/s. Để đi qua cầu, tàu phải đi quãng đường bằng chiều dài cầu cộng chiều dài tàu: $480+180=660$ (m), mất $\dfrac{660}{16}=41{,}25$ (s) $\ne30$ (s). Sai.}
\end{ex}
